\origpage{44}{ROBERTSON AND SEYMOUR}
\textit{Subproof.} If there exists $v \in V(H) \cap Z$ incident with $r$, then $d(v, r) \le 1$, and so
\begin{equation*}
d(z, r) \le d(z, v) + d(v, r) \le (\kappa - 1) + 1 = \kappa.
\end{equation*}
If there is no such $v$, then since $\bar{r} \cap Z \ne \emptyset$, it follows that $\bar{r} \cap z' \ne \emptyset$, and hence $d(z, r) \le 2 \le \kappa$. Thus in either case $d(z, r) \le \kappa \le \theta - 3$, and so $r \subseteq \Delta - bd(\Delta)$ by (1). This proves (2).

From (2) and (5.2) we deduce that there is a circuit $C$ of $H$ with $U(C) \subseteq \Delta$, bounding an open disc $\Lambda \subseteq \Delta$ including $Z$, and such that every edge of $C$ is incident with a region $r$ of $H$ satisfying $\bar{r} \cap Z \ne \emptyset$. Moreover, $x \subseteq \Lambda$ for every $x \in A(H)$ with $d(z, x) < \kappa$; for if $x'$ is the corresponding atom of $K$, then $x' \subseteq Z \subseteq \Lambda$, and so $x \cap \Lambda \ne \emptyset$, and therefore $x \subseteq \Lambda$, as claimed. It remains to show that $\Lambda$ is a $(\kappa + 2)$-zone around $z$. To see this, let $Y$ be the union of $ins(W)$ taken over all closed walks $W$ of length $< 2(\kappa + 3)$ with $z' \subseteq ins(W)$. We claim that $U(C) \subseteq Y$. Since $Y$ is closed, it suffices to show that $e \subseteq Y$ for all $e \in E(C)$. Let $e \in E(C)$, and let $r$ be the region of $H$ with $r \subseteq \Lambda$, incident with $e$. Then, by choice of $C$, $\bar{r} \cap Z \ne \emptyset$ and so $d(z, r) \le \kappa$ by (2). Since $d(r, e) \le 2$ it follows that $d(z, e) \le \kappa + 2$. Hence $Y$ includes the region of $K$ including $e$, and hence includes $e$. This proves that $U(C) \subseteq Y$. Since $Y$ is simply-connected by [5, theorem (8.10)], there is a closed disc $\Delta^* \subseteq Y$ with $bd(\Delta^*) = U(C)$. Since $\Delta^* \subseteq Y \subseteq \Delta$ and $\bar{\Lambda} \subseteq \Delta$, and the closed discs $\Delta^*$, $\bar{\Lambda} \subseteq \Delta$ are both bounded by $U(C)$, it follows that $\Delta^* = \bar{\Lambda}$, and so $\bar{\Lambda} \subseteq Y$. Now let $x \in A(H)$ with $x \cap \bar{\Lambda} \ne \emptyset$. Then $x \cap Y \ne \emptyset$, and so $d(z, x) \le \kappa + 2$. Hence $\Lambda$ is a $(\kappa + 2)$-zone as required. $\square$

The next result is not needed for this paper, and is included just for future reference.

\stmt
(9.3) \textit{Let $H$ be a 2-cell rigid drawing in $\Sigma$, and let $\mathscr{T}$ be a respectful tangle in $H$ of order $\theta$, with metric $d$. Let $z \in A(H)$, and let $\kappa$ be an integer with $2 \le \kappa \le \theta - 6$. Then there is a $(\kappa + 5)$-zone $\Lambda$ around $z$ such that $H \cap (\Sigma - \Lambda)$ is rigid and $x \subseteq \Lambda$ for every $x \in A(H)$ with $d(z, x) < \kappa$.}

\textit{Proof.} By two applications of (9.2) (with $\kappa$ and $\kappa + 3$) there is a $(\kappa + 2)$-zone $\Lambda_1$ around $z$ such that $x \subseteq \Lambda_1$ for every $x \in A(H)$ with $d(z, x) < \kappa$, and there is a $(\kappa + 5)$-zone $\Lambda_2$ around $z$ such that $x \subseteq \Lambda_2$ for every $x \in A(H)$ with $d(z, x) < \kappa + 3$. Since $\Lambda_1$ is a $(\kappa + 2)$-zone it follows that $\bar{\Lambda}_1 \subseteq \Lambda_2$. Choose a circuit $C$ of $H$ with $U(C) \subseteq \Lambda_2$, bounding an open disc $\Lambda \subseteq \Lambda_2$ with $\Lambda_1 \subseteq \Lambda$, and such that $\Lambda$ is maximal. Then $\Lambda$ is a $(\kappa + 5)$-zone around $z$ since $\Lambda \subseteq \Lambda_2$, and $x \subseteq \Lambda_1 \subseteq \Lambda$ for every $x \in A(H)$ with $d(z, x) < \kappa$, and it remains to show that $H \cap (\Sigma - \Lambda)$ is rigid. Suppose not; then by (6.3) there is a region $r$ of $H$ with $r \cap \Lambda = \emptyset$ such that $\bar{r} \cap \bar{\Lambda} \ne \emptyset$
